\section{Higher-power norm classes and finite arithmetic descent}
\label{pl-section}

The second actual mixed norm is the binary quartic
\[
 F(a,b)=a^4+3a^3b+5a^2b^2+3ab^3+b^4,
 \qquad F(x)=F(x,1).
\]
Its discriminant is $117$. We first give the exact denominator correction
for a direct affine lift. We then use integral primitivity and finite
Kummer descent to remove that denominator restriction from the
fixed-exponent finiteness theorem. This is an explicit application of
the classical Darmon--Granville mechanism
\cite[Theorem~1 and its consequence on p.~514]{DarmonGranville1995},
not a new general finiteness theorem. Throughout, a seed means positive
coprime integers $a,b$; its height is $a+b$.

\subsection{The exact direct lift}

For $h\ge2$ put
\[
 \delta=\gcd(h,4),\qquad m=h/\delta,\qquad k=4/\delta.
\]
A positive integer is \emph{$h$-free} when every prime valuation belongs
to $\{0,\ldots,h-1\}$. Every positive rational class modulo $h$-th powers
has a unique positive $h$-free integer representative. For $m=1$, the
only $m$-free representative is $1$.

\begin{lemma}[Integral lifting and the denominator class]\label{pl-direct-lift}
Fix an $h$-free positive integer $D$ and an $m$-free positive integer
$\beta$. Primitive positive solutions
\begin{equation}\label{pl-integer-lift}
 F(a,b)=DW^h,\qquad b=\beta B^m,\qquad W,B\in\mathbb Z_{>0}
\end{equation}
correspond bijectively to positive rational points $(x,y)$ satisfying
\begin{equation}\label{pl-affine-lift}
 y^h=\frac{\beta^4F(x)}D
\end{equation}
whose reduced denominator $b$ has the indicated $m$-th power class.
The maps are $x=a/b$, $y=W/B^k$, and conversely $W=yB^k$.
The denominator-class map
\[
 \mathbb Q_{>0}^{\times}/\mathbb Q_{>0}^{\times m}
 \longrightarrow
 \mathbb Q_{>0}^{\times}/\mathbb Q_{>0}^{\times h},
 \qquad [b]\longmapsto[b^4]
\]
is injective.
\end{lemma}
\begin{proof}
The identity $4m=hk$ gives $b^4=\beta^4(B^k)^h$; homogeneity proves
the displayed maps. To check integral lifting, suppose $q>0$ is rational,
$Dq^h$ is an integer, and $D$ is $h$-free. For every prime $p$,
\[
 h v_p(q)=v_p(Dq^h)-v_p(D)\ge-(h-1).
\]
The left side is a multiple of $h$, hence is nonnegative. Thus $q$ is
an integer. Apply this to $B$ and then to $W$; when $m=1$ the first
step is immediate. Positivity and the unique reduced fraction $a/b$
give the bijection. Finally, $h\mid4e$ if and only if $m\mid e$ for
each integer $e$, proving the class-map assertion prime by prime.
\end{proof}

The correction is necessary. The primitive seed $(a,b)=(1,2)$ satisfies
$F(a,b)=67$. With $D=67$ and $W=1$,
\[
 \frac{F(1/2)}{67}=\frac1{16},
\]
which is not a rational $h$-th power when $h\nmid4$. Its $2$-adic
valuation is $-4$. Thus $W/b^{4/h}$ cannot be used as a rational
coordinate in general. This is a counterexample to that direct map,
not to the integer finiteness theorem below.

\begin{theorem}[Direct-cover genus and coupled class]\label{pl-direct-genus}
For $d>0$ rational, the smooth projective curve with function field
$\mathbb Q(x)(y)$, $y^h=F(x)/d$, is geometrically connected and has genus
\begin{equation}\label{pl-direct-genus-formula}
 g_h=\frac{3h-2-\gcd(h,4)}2.
\end{equation}
For fixed $h\ge3$, only finitely many primitive positive seeds
$F(a,b)=VW^h$ can have the coupled class
$[Vb^{-4}]_h$ in any fixed finite set of $h$-th power classes.
In particular fixing $[V]_h$ and $[b]_m$ suffices.
\end{theorem}
\begin{proof}
Over the algebraic closure, valuation one at a root gives Kummer degree
$h$. The four simple roots have inertia $h$, contributing $4(h-1)$.
At infinity the pole order is four: there are $\delta$ points of
ramification index $h/\delta$, contributing $h-\delta$. There is no
other ramification. Thus
\[
 2g_h-2=-2h+4(h-1)+(h-\delta)=3h-4-\delta.
\]
For $h\ge3$, this genus is at least three, so Faltings's theorem gives
finitely many rational points on each fixed twist. If
$Vb^{-4}=dq^h$, then $F(a/b)=d(qW)^h$, which supplies such a rational
point. Each positive reduced rational $a/b$ determines its seed uniquely.
A finite union over twists proves the result.
\end{proof}

For $h=4$, $m=1$, so the denominator correction is automatic and each
fixed fourth-power class of $M_1$ gives only finitely many seeds. For
$h=2$, the genus is one, in agreement with the infinite square family
of Section~\ref{qg-section}. The next theorem strengthens the
fixed-class conclusion for every $h\ge3$ without fixing a denominator
class.

\subsection{Finite descent using the four linear factors}

\begin{lemma}[Finite Kummer classes]\label{pl-selmer-finite}
Let $K$ be a number field, $S$ a finite set of its nonarchimedean primes,
and $h\ge2$. Then
\[
 K(S,h)=
 \{z\in K^\times:v_{\mathfrak p}(z)\in h\mathbb Z
       \text{ for }\mathfrak p\notin S\}/K^{\times h}
\]
is finite.
\end{lemma}
\begin{proof}
In the ring $\mathcal O_{K,S}$, write $(z)=\mathfrak a^h$. The map
to $[\mathfrak a]$ has image in the finite group
$\operatorname{Cl}(\mathcal O_{K,S})[h]$. Its kernel is the group of
$S$-units modulo $h$-th powers: if $\mathfrak a=(t)$, then $z/t^h$
is an $S$-unit, and conversely an $S$-unit has trivial ideal. The
$S$-unit group is finitely generated. Thus both kernel and image are
finite. This uses finiteness of the ideal class group, with no
principal-ideal or element-factorization assumption.
\end{proof}

\begin{theorem}[Finite seeds for a fixed higher-power residual class]
\label{pl-fixed-residual}
Fix $h\ge3$ and a positive rational $h$-th power class $[D]_h$.
There are only finitely many primitive positive seeds satisfying
\begin{equation}\label{pl-fixed-equation}
 F(a,b)=DW^h,\qquad W\in\mathbb Q_{>0}.
\end{equation}
The conclusion holds for a fixed finite set of residual classes,
and also when the positive integer residual has prime support in
one fixed finite set. No denominator-class hypothesis is required.
\end{theorem}
\begin{proof}
Choose the positive $h$-free integer representative $D$; the integral
lifting argument above makes $W$ an integer. Let $K$ contain the four
roots $\alpha_1,\ldots,\alpha_4$ of $F$ and the $h$-th roots of unity.
These roots are integral because $F$ is monic. Put
\[
 L_i=a-\alpha_i b,\qquad \prod_{i=1}^4L_i=DW^h.
\]
Let $S$ consist of the primes of $K$ over rational primes dividing
$117D$. Outside $S$ all differences $\alpha_i-\alpha_j$ are units,
since their squared product is $117$. If a prime $\mathfrak p\notin S$
divided both $L_i$ and $L_j$, it would divide $b$ and then $a$,
contradicting an integer Bezout identity for the primitive seed.
Hence at most one $L_i$ has positive valuation at $\mathfrak p$.
Taking valuations of the product shows
$h\mid v_{\mathfrak p}(L_i)$ for every $i$ outside $S$.

Lemma~\ref{pl-selmer-finite} gives only finitely many possible classes
of the $L_i$, and hence only finitely many triples of classes of
$L_i/L_4$ for $i=1,2,3$. Choose their representatives $\xi_i$.
Every seed gives a $K$-point, with $x=a/b$, on one of the finitely many
curves
\begin{equation}\label{pl-four-branch-cover}
 y_i^h=\frac{x-\alpha_i}{\xi_i(x-\alpha_4)},
 \qquad i=1,2,3.
\end{equation}
These points avoid the branch values since $F(x)>0$ for $x>0$.

Over an algebraic closure the three Kummer classes are independent
modulo $h$: evaluating the valuation of a relation at $\alpha_i$
forces its $i$-th exponent to be divisible by $h$. This holds also
for composite $h$. The smooth projective normalization of
\eqref{pl-four-branch-cover} is therefore geometrically connected of
degree $h^3$, with group $(\mathbb Z/h\mathbb Z)^3$.

At each $\alpha_i$ for $i\le3$ the inertia order is $h$. At
$\alpha_4$ the three functions all have a simple pole. Geometric
local units have $h$-th roots in characteristic zero, so the three
extensions there adjoin the same $h$-th root of a uniformizer.
The inertia group is diagonal cyclic of order $h$, not of order
$h^3$. At infinity all three ratios are units, and there is no
ramification. Consequently the four branch points each contribute
$h^2(h-1)$, and Riemann--Hurwitz gives
\begin{equation}\label{pl-descent-genus}
 2G_h-2=-2h^3+4h^2(h-1),\qquad G_h=1+h^2(h-2)>1.
\end{equation}
Faltings's theorem over the fixed number field $K$ gives finitely many
$K$-points on each curve. There are finitely many curves, and each
positive reduced rational $x$ determines $a,b$ uniquely. This proves
the fixed-class statement; a finite union proves its finite-class
version. If the residual support is in $S_0$, its $h$-free
representative has at most $h^{|S_0|}$ possibilities, proving the
last assertion.
\end{proof}

The arithmetic covers capture every actual primitive seed. They do
not assert that the direct coordinate $W/b^{4/h}$ is rational.
Primitivity outside a finite set supplies the extra information that
the direct affine normalization lacked.

\subsection{An exponent restriction for possible unbounded families}

\begin{corollary}[Residual support or exponent escape]\label{pl-exponent-escape}
Consider actual second profiles $M_1=V_1Q_1^{g_1}$ along a sequence
of primitive positive seeds with $a+b\longrightarrow\infty$.
For every fixed $h\ge3$, among terms with $h\mid g_1$, the class
$[V_1]_h$ eventually leaves every fixed finite set of classes.
If all $V_1$ have prime support in one fixed finite set, then each
fixed $h\ge3$ divides $g_1$ only finitely often. In particular,
eventually $v_2(g_1)\le1$, and for each fixed $B$ no odd prime
$p\le B$ divides $g_1$ at sufficiently large indices. If also
$g_1\longrightarrow\infty$, its least odd prime factor tends to infinity.
\end{corollary}
\begin{proof}
For $h\mid g_1$, substitute $W=Q_1^{g_1/h}$ in
Theorem~\ref{pl-fixed-residual}. Each indicated finite class set gives
only finitely many seeds, which a height-tending-to-infinity sequence
eventually leaves. With fixed residual support, apply this to $h=4$
and each of the finitely many odd primes at most $B$. If $g_1$ tends
to infinity, after excluding multiples of four the values with no odd
prime divisor, namely $1$ and $2$, no longer occur.
\end{proof}

\begin{corollary}[Finite total support]\label{pl-total-support}
For each fixed finite prime set $S$, there are only finitely many
primitive positive seeds with $\operatorname{supp}(M_1)\subseteq S$.
Thus, along a sequence with $a+b\longrightarrow\infty$, if
$\operatorname{supp}(Q_1)\subseteq S_Q$ for one fixed finite set,
then for every fixed finite $S_V$ one eventually has
$\operatorname{supp}(V_1)\not\subseteq S_V$. No condition on
divisibility of $g_1$ is needed.
\end{corollary}
\begin{proof}
Write $M_1=DW^3$ with $D$ cube-free. There are at most $3^{|S|}$
possible $D$, so Theorem~\ref{pl-fixed-residual} applies. For the
profile consequence use $S=S_Q\cup S_V$; a height-tending-to-infinity
sequence eventually leaves the resulting finite set of seeds.
\end{proof}

This is a concrete Thue--Mahler consequence, with no quantitative
prime-growth estimate. The unit residual case $V_1=1$ is included in
Corollary~\ref{pl-exponent-escape}, even with arbitrary moving
roots. Each fixed exponent divisor $h\ge3$ supports only finitely many
seeds. This does not assert that such seeds do not exist, and it is
not uniform in $h$. Large odd prime exponents, twice such exponents,
and moving residual support remain outside this obstruction. A small
$\lambda_1=\max\{1,\log V_1\}/g_1$ need not bound the residual support
or its power class, so the two-step finite threshold remains unresolved.

The number-field and geometric dependencies here are the finite ideal
class group, finite generation of $S$-units, Kummer theory,
Riemann--Hurwitz, and Faltings's theorem. This section supplies ordinary
proofs; it does not claim Lean verification of those dependencies or
of the resulting finiteness theorem. The companion exact arithmetic
replay checks the denominator identities and examples only.
